\scanpage{31}
% archive sheet number: 30
\mspage{\handwritten{26}}
% L1
Démonstration.

% L2–8
a) et b). La démonstration de a) et de la première assertion dans b) est essentiellement identique à celle de 1.7. a) et b), et nous nous dispensons de reproduire ici le raisonnement. Signalons seulement qu'il faut faire appel dans a) à Exp IX 5.6. (utilisant le fait que $G$ est à fibres connexes). Prouvons la deuxième assertion de b), i.e. que si $Z$ est un centre réductif de $G$, alors $G'=G/Z$ admet pour centre réductif le sous-groupe unité.
% L8–15
Notons tout de suite qu'en vertu de Exp VIII 5.1. le groupe quotient $G/Z$ estb bien représentable, il est affine sur $S$, de présentation finie sur $S$ (Exp VIII 5.8.) et même lisse sur $S$ (car il suffit de le voir pour les fibres, $G'$ étant plat de présentation finie sur $S$, or pour les fibres \deleted{cela est} on a signalé dans Exp VI qu'un quotient d'un groupe algébrique lisse sur un corps $k$ est lisse); enfin $G$ étant à fibres connexes, il en est de même de $G'$. Ainsi $G'$ satisfait aux mêmes hypothèses de départ que $G$.
% L15–21
Pour voir que le sous-groupe unité de $G'$ est un centre réductif, on peut se borner si on veut au cas où $\deleted{\unsure{s}}\handwritten{S}$ est le spectre d'un corps (4.\deleted{\unsure{2|3}}\handwritten{3}.). Soit $Z'$ un sous-groupe central de type multiplicatif de $G'$, tout revient à prouver que $\handwritten{Z}'$ est réduit au groupe unité. Soit $Z_{1}$ \deleted{\unsure{l'im}} l'image inverse de $Z'$ dans $G$, et considérons les opérations de $Z_{1}$ sur $G$ induites par automorphismes intérieurs de $G$. Comme $Z$ est central, $Z$ opère \deleted{\unsure{tri}} trivialement, de sorte que $Z_{1}$ opère par l'intermédiaire d'opérations de $Z'$ sur $G$.
% L22–26
D'ailleurs $Z$ étant dans le centre de $G$, $Z\deleted{\unsure{'}}_{1}$ donc $Z'$ opère trivialement sur $Z$, de plus $Z'$ étant dans le centre de $G'$, les opérations de $Z'$ dans le quotient $G/Z=G'$ \handwritten{\unsure{/|(}}sont également triviales. Comme $Z$ est central, il s'ensuit aussitôt que les opérations de $Z'$ sur $G$ proviennent d'un
